\documentclass{stex}

\libusepackage{naproche}
\libinput{preamble}

\begin{document}
\begin{smodule}{zfc.ftl}

\importmodule[libraries/set-theory]{axioms?set-existence.ftl}
\importmodule[libraries/set-theory]{axioms?separation.ftl}
\importmodule[libraries/set-theory]{axioms?pairing.ftl}
\importmodule[libraries/set-theory]{axioms?union.ftl}
\importmodule[libraries/set-theory]{axioms?powerset.ftl}
\importmodule[libraries/set-theory]{axioms?infinity.ftl}
\importmodule[libraries/set-theory]{axioms?replacement.ftl}
\importmodule[libraries/set-theory]{axioms?foundation.ftl}
\importmodule[libraries/set-theory]{axioms?choice.ftl}
\importmodule[libraries/foundations]{pairs-and-products.ftl}
\importmodule[libraries/foundations]{invertible-maps.ftl}

\begin{proposition}[forthel]
  $\EmptyClass$ is a set.
\end{proposition}
\begin{proof}[forthel]
  Take a set $x$ (by \sn{Set Existence Axiom}).
  Define $A \DefEq \SClass{y}{x}{y \NotEq y}$.
  Then $A$ is a set (by \sn{Separation Axiom}).
  We have $A \Eq \EmptyClass$.
  Hence $\EmptyClass$ is a set.
\end{proof}

\begin{proposition}[forthel]
  Let $x, y$ be sets.
  Then $x \Union y$ is a set.
\end{proposition}
\begin{proof}[forthel]
  Take $X \Eq \Pair{x}{y}$.
  Then $X$ is a set.
  Hence $\UnionOver X$ is a set (by \sn{Union Axiom}).
  Indeed $X$ is a system of sets.
  We have $x \Union y \Eq \UnionOver X$.
  Thus $x \Union y$ is a set.
\end{proof}

\begin{proposition}[forthel]
  Let $x, y$ be sets.
  Then $x \Intersect y$ is a set.
\end{proposition}
\begin{proof}[forthel]
  We have $x \Intersect y \Subclass x$.
  Hence $x \Intersect y$ is a set (by \sn{Separation Axiom}).
\end{proof}

\begin{proposition}[forthel]
  Let $x, y$ be sets.
  Then $x \Diff y$ is a set.
\end{proposition}
\begin{proof}[forthel]
  We have $x \Diff y \Subclass x$.
  Hence $x \Diff y$ is a set (by \sn{Separation Axiom}).
\end{proof}

\begin{proposition}[forthel]
  Let $x, y$ be sets.
  Then $x \Prod y$ is a set.
\end{proposition}
\begin{proof}[forthel]
  $\Singleton{a}$ and $\Pair{a}{b}$ are sets for each $a \Elem x$ and each $b \Elem y$.
  Define $P \DefEq \CClass{\Pair{\Singleton{a}}{\Pair{a}{b}}}{a \Elem x\text{ and }b \Elem y}$.

  (1) $P$ is a set.
  \begin{proof}
    Let us show that $P \Subclass \Pow(\Pow(x \Union y))$.
      Let $p \Elem P$.
      Consider $a \Elem x$ and $b \Elem y$ such that $p \Eq \Pair{\Singleton{a}}{\Pair{a}{b}}$.
      Then $a, b \Elem x \Union y$.
      Hence $\Singleton{a}, \Pair{a}{b} \Elem \Pow(x \Union y)$.
      Thus $\Pair{\Singleton{a}}{\Pair{a}{b}} \Elem \Pow(\Pow(x \Union y))$.
    End.

    $x \Union y$ is a set.
    Consequently $\Pow(\Pow(x \Union y))$ is a set (by \sn{Powerset Axiom}).
    Therefore $P$ is a set (by \sn{Separation Axiom}).
  \end{proof}

  Define $l(p) \DefEq$ ``choose $a \Elem x$, choose $b \Elem y$ such that $p \Eq \Pair{\Singleton{a}}{\Pair{a}{b}}$ in $a$'' for $p \Elem P$.
  Define $r(p) \DefEq$ ``choose $a \Elem x$, choose $b \Elem y$ such that $p \Eq \Pair{\Singleton{a}}{\Pair{a}{b}}$ in $b$'' for $p \Elem P$.

  Define $f(p) \DefEq (l(p), r(p))$ for $p \Elem P$.

  Let us show that for any objects $u, u', v, v'$ if
  $\Pair{\Singleton{u}}{\Pair{u}{v}} \Eq \Pair{\Singleton{u'}}{\Pair{u'}{v'}}$ then $u \Eq u'$ and $v \Eq v'$.
    Let $u, u', v, v'$ be objects.
    Assume $\Pair{\Singleton{u}}{\Pair{u}{v}} \Eq \Pair{\Singleton{u'}}{\Pair{u'}{v'}}$.
    Then ($\Singleton{u} \Eq \Singleton{u'}$ or $\Singleton{u} \Eq \Pair{u'}{v'}$) and ($\Pair{u}{v} \Eq \Singleton{u'}$ or $\Pair{u}{v} \Eq \Pair{u'}{v'}$).
    Thus ($\Singleton{u} \Eq \Singleton{u'}$ and ($\Pair{u}{v} \Eq \Singleton{u'}$ or $\Pair{u}{v} \Eq \Pair{u'}{v'}$)) or ($\Singleton{u} \Eq \Pair{u'}{v'}$ and ($\Pair{u}{v} \Eq \Singleton{u'}$ or $\Pair{u}{v} \Eq \Pair{u'}{v'}$)).

    \begin{case}{$\Singleton{u} \Eq \Singleton{u'}$ and ($\Pair{u}{v} \Eq \Singleton{u'}$ or $\Pair{u}{v} \Eq \Pair{u'}{v'}$).}
      We have $\Singleton{u} \Eq \Singleton{u'}$.
      Hence $u \Eq u'$.

      \begin{case}{$\Pair{u}{v} \Eq \Singleton{u'}$.}
        Then $u \Eq u'\Eq v$.
        Hence $\Pair{\Singleton{u}}{\Pair{u}{u}} \Eq \Pair{\Singleton{u}}{\Pair{u}{v'}}$ (by 1).
        Thus $\Singleton{\Singleton{u}} \Eq \Pair{\Singleton{u}}{\Pair{u}{v'}}$.
        Therefore $\Singleton{u} \Eq \Pair{u}{v'}$.
        Consequently $v' \Eq u \Eq v$.
      \end{case}

      \begin{case}{$\Pair{u}{v} \Eq \Pair{u'}{v'}$.}
        Then $\Pair{u}{v} \Eq \Pair{u}{v'}$.
        Hence $v \Eq v'$.
      \end{case}
    \end{case}

    \begin{case}{$\Singleton{u} \Eq \Pair{u'}{v'}$ and ($\Pair{u}{v} \Eq \Singleton{u'}$ or $\Pair{u}{v} \Eq \Pair{u'}{v'}$).}
      We have $\Singleton{u} \Eq \Pair{u'}{v'}$.
      Hence $u \Eq u'$.

      \begin{case}{$\Pair{u}{v} \Eq \Singleton{u'}$.}
        Then $u \Eq v \Eq u'$.
        Hence $v \Eq v'$.
      \end{case}

      \begin{case}{$\Pair{u}{v} \Eq \Pair{u'}{v'}$.}
        Then $\Pair{u}{v} \Eq \Pair{u}{v'}$.
        Hence $v \Eq v'$.
      \end{case}
    \end{case}
  End.

  $\Pair{\Singleton{a}}{\Pair{a}{b}} \Elem \Dom(f)$ for any $a \Elem x$ and any $b \Elem y$. % Needed for ontological checking
  \begin{proof}
    Let $a \Elem x$ and $b \Elem y$.
    Then $\Pair{\Singleton{a}}{\Pair{a}{b}} \Elem P$.
  \end{proof}

  Let us show that for any $a \Elem x$ and any $b \Elem y$ we have $f(\Pair{\Singleton{a}}{\Pair{a}{b}}) \Eq (a,b)$.
    Let $a \Elem x$ and $b \Elem y$.
    Take $p \Eq \Pair{\Singleton{a}}{\Pair{a}{b}}$.
    Then $p$ is a set.
    Then we can choose $a' \Elem x$ and $b' \Elem y$ such that $p \Eq \Pair{\Singleton{a'}}{\Pair{a'}{b'}}$ and $l(p) \Eq a'$.
    Then $a \Eq a'$ and $b \Eq b'$.
    Hence $l(p) \Eq a$.
    Choose $a'' \Elem x$ and $b'' \Elem y$ such that $p \Eq \Pair{\Singleton{a''}}{\Pair{a''}{b''}}$ and $r(p) \Eq b''$.
    Then $a \Eq a''$ and $b \Eq b''$.
    Thus $r(p) \Eq b$.
    Therefore $f(p) \Eq (a,b)$.
  End.

  (2) $x \Prod y \Eq \Im{f}{P}$.
  \begin{proof}
    For all $p \Elem P$ we have $l(p) \Elem x$ and $r(p) \Elem y$.
    Hence $f(p) \Elem x \Prod y$ for all $p \Elem P$.
    Therefore $\Im{f}{P} \Subclass x \Prod y$.

    Let us show that $x \Prod y \Subclass \Im{f}{P}$.
      Let $z \Elem x \Prod y$.
      Take $a \Elem x$ and $b \Elem y$ such that $z \Eq (a,b)$.
      Then $(a,b) \Eq f(\Pair{\Singleton{a}}{\Pair{a}{b}})$.
      Hence there exists a $p \Elem P$ such that $(a,b) \Eq f(p)$.
      Thus $(a,b) \Elem \Im{f}{P}$.
    End.

    Consequently $x \Prod y \Eq \Im{f}{P}$.
  \end{proof}

  Thus $x \Prod y$ is the image of some set under some map.
  Therefore $x \Prod y$ is a set (by \sn{Replacement Axiom}).
\end{proof}

\begin{proposition}[forthel]
  Let $X$ be a nonempty system of sets.
  Then $\IntersectOver X$ is a set.
\end{proposition}
\begin{proof}[forthel]
  Take an element $x$ of $X$.
  Then $\IntersectOver X \Subclass x$.
  Hence $\IntersectOver X$ is a set (by \sn{Separation Axiom}).
\end{proof}

\begin{proposition}[forthel]
  Let $f$ be a map such that $\Dom(f)$ is a set.
  Then $\Range(f)$ is a set.
\end{proposition}
\begin{proof}[forthel]
  $\Range(f) \Eq \Im{f}{\Dom(f)}$ and $\Im{f}{\Dom(f)}$ is a set.
  Hence $\Range(f)$ is a set (by \sn{Replacement Axiom}).
\end{proof}

\begin{proposition}[forthel]
  Let $A$ be a class and $x$ be a set.
  Assume that there exists an injective map from $A$ to $x$.
  Then $A$ is a set.
\end{proposition}
\begin{proof}[forthel]
  Consider an injective map $f$ from $A$ to $x$.
  Then $\Inv{f}$ is a bijection between $\Range(f)$ and $A$.
  $\Range(f)$ is a set and $A$ is the image of $\Range(f)$ under $\Inv{f}$.
  Indeed $\Range(f) \Subclass x$.
  Thus $A$ is a set (by \sn{Replacement Axiom}).
\end{proof}

\begin{proposition}[forthel]
  There exist no sets $x, y$ such that $x \Elem y$ and $y \Elem x$.
\end{proposition}
\begin{proof}[forthel]
  Assume the contrary.
  Take sets $x,y$ such that $x \Elem y$ and $y \Elem x$.
  Consider an element $z$ of $\Pair{x}{y}$ such that $\Pair{x}{y}$ and $z$ are disjoint (by \sn{Foundation Axiom}).
  Indeed $\Pair{x}{y}$ is a nonempty system of sets.
  Then we have $z \Eq x$ or $z \Eq y$.

  \begin{case}{$z \Eq x$.}
    Then $x$ and $\Pair{x}{y}$ are disjoint.
    Hence $y \NotElem x$.
    Contradiction.
  \end{case}

  \begin{case}{$z \Eq y$.}
    Then $y$ and $\Pair{x}{y}$ are disjoint.
    Hence $x \NotElem y$.
    Contradiction.
  \end{case}
\end{proof}

\begin{corollary}[forthel]
  Let $x$ be a set.
  Then $x \NotElem x$.
\end{corollary}

\begin{proposition}[forthel]
  Let $f, g$ be functions.
  Assume that $\Dom(f) \Eq \Dom(g)$ and $f(a) \Eq g(a)$ for all $a \Elem \Dom(f)$.
  Then $f \Eq g$.
\end{proposition}

\begin{proposition}[forthel]
  Let $x, y$ be sets.
  Then $\FunSpace{x}{y}$ is a set.
\end{proposition}
\begin{proof}[forthel]
  Define $R \DefEq \SClass{F}{\Pow(x \Prod y)}{\text{ (for all }a \Elem x\text{ there exists a }b \Elem y\text{ such that }(a,b) \Elem F\text{) and for all }a \Elem x\text{ and all }b, b' \Elem y\text{ such that }(a,b), (a,b') \Elem F\text{ we have }b \Eq b'}$.

  [prover vampire]
  Every element of $R$ is a set. % Needed for ontological checking
  Define $h(F) \DefEq \Map{a}{x}$ ``choose $b \Elem y$ such that $(a,b) \Elem F$ in $b$'' for $F \Elem R$.
  [prover eprover]

  Let us show that $\FunSpace{x}{y} \Subclass \Range(h)$.
    Let $f \Elem \FunSpace{x}{y}$.
    Define $F \DefEq \CClass{(a,f(a))}{a \Elem x}$.

    Then $F \Elem R$.
    \begin{proof}
      Define $g(a) \DefEq (a,f(a))$ for $a \Elem x$.
      Then $F \Eq \Range(g)$.
      Hence $F$ is a set.
      Thus $F \Elem \Pow(x \Prod y)$.
      Indeed $F \Subclass x \Prod y$.

      (1) For all $a \Elem x$ there exists a $b \Elem y$ such that $(a,b) \Elem F$.

      (2) For all $a \Elem x$ and all $b, b' \Elem y$ such that $(a,b),
      (a,b') \Elem F$ we have $b \Eq b'$.
      \begin{proof}
        Let $a \Elem x$ and $b, b' \Elem y$.
        Assume $(a,b), (a,b') \Elem F$.
        Take $c, c' \Elem F$ such that $c \Eq (a,b)$ and $c' \Eq (a,b')$.
        Take $d, d' \Elem x$ such that $c \Eq (d,f(d))$ and $c' \Eq (d',f(d'))$.
        Then $a \Eq d$ and $b \Eq f(d)$ and $a \Eq d'$ and $b' \Eq f(d)$.
        Hence $b \Eq b'$.
      \end{proof}

      [prover vampire]
      Hence the thesis.
    \end{proof}

    $\Dom(f) \Eq \Dom(h(F))$ and for each $a \Elem \Dom(f)$ we have $h(F)(a) \Eq f(a)$.
    Hence $f \Eq h(F)$.
    Thus $f \Elem \Range(h)$.
  End.

  Therefore $\FunSpace{x}{y}$ is a set.
  Indeed $R$ is a set.
\end{proof}

\end{smodule}
\end{document}
